% ch22: 定积分概念
\chapter{定积分的概念——从面积到黎曼和}
\begin{applicationbox}
学完本章：理解黎曼和与定积分定义，掌握几何意义
\end{applicationbox}

\section{面积问题} 求曲线下的面积——用矩形逼近：\(\sum f(x_i^*)\Delta x_i\)，当分割无限细时取极限。

\section{定积分定义} \(\int_a^b f(x)dx = \lim_{\max\Delta x_i\to0}\sum_{i=1}^n f(\xi_i)\Delta x_i\)

\textbf{用 ε-δ 语言：} 对任意 \(\varepsilon>0\)，存在 \(\delta>0\)，
使得当分割的最大宽度 \(\|\Delta\|<\delta\) 时，对任意取点方式 \(\xi_i\)，
\[
\left|\sum_{i=1}^n f(\xi_i)\Delta x_i - \int_a^b f(x)dx\right| < \varepsilon
\]

\section{连续函数必可积（ε-δ 证明）}
设 \(f\) 在 \([a,b]\) 上连续，则 \(f\) 一致连续：
对任意 \(\varepsilon>0\)，存在 \(\delta>0\)，使 \(|x-y|<\delta\) 时 \(|f(x)-f(y)|<\frac{\varepsilon}{b-a}\)。
取分割 \(a=x_0<x_1<\cdots<x_n=b\) 满足最大宽度 \(<\delta\)。记
\(M_i=\max_{[x_{i-1},x_i]}f\)，\(m_i=\min_{[x_{i-1},x_i]}f\)，则
\[
U-L = \sum (M_i-m_i)\Delta x_i < \frac{\varepsilon}{b-a}\sum\Delta x_i = \varepsilon
\]
由达布定理，\(f\) 在 \([a,b]\) 上黎曼可积。\(\blacksquare\)

\section{几何意义} \(\int_a^b f(x)dx\) = 曲线与 \(x\) 轴围成的面积（若 \(f\ge0\)）

\begin{examplebox}[1——黎曼和近似]
\(\int_0^1 x^2 dx\) 用 \(n=4\) 个右端点矩形：\(\frac14[(0.25)^2+(0.5)^2+(0.75)^2+1^2]=0.46875\)，精确值 \(1/3\approx0.333\)。
\end{examplebox}

\section{习题}
\begin{exercise}[0] 定积分 \(\int_a^b f(x)dx\) 的几何意义是什么？\end{exercise}
\begin{exercise}[0] 黎曼和中 \(\Delta x_i\) 代表什么？\end{exercise}
\begin{exercise}[0] 若 \(f(x)\ge0\)，\(\int_a^b f(x)dx\) 的符号是？\end{exercise}
\begin{exercise}[0] 定积分 \(\int_a^a f(x)dx\) 等于多少？\end{exercise}
\begin{exercise}[1] 用右端点矩形近似 \(\int_0^1 x dx\)（\(n=4\)）。\end{exercise}
\begin{exercise}[1] 用左端点矩形近似 \(\int_0^1 x^2 dx\)（\(n=4\)）。\end{exercise}
\begin{exercise}[1] 求 \(\int_0^2 3 dx\) 的精确值（几何意义）。\end{exercise}
\begin{exercise}[1] 求 \(\int_0^1 x dx\) 的精确值。\end{exercise}
\begin{exercise}[2] 用 \(n=6\) 的黎曼和近似 \(\int_0^\pi \sin x dx\)。\end{exercise}
\begin{exercise}[2] 求 \(\int_0^a x dx\) 的值。\end{exercise}
\begin{exercise}[2] 求 \(\int_{-a}^a x^3 dx\) 的值（利用奇函数性质）。\end{exercise}
\begin{exercise}[2] 用几何意义求 \(\int_{-1}^1 \sqrt{1-x^2} dx\)。\end{exercise}
\begin{exercise}[3] 用定义计算 \(\int_0^1 x dx\)（取等分分割、右端点）。\end{exercise}
\begin{exercise}[3] 用定义计算 \(\int_0^1 x^2 dx\)。\end{exercise}
\begin{exercise}[3] 比较 \(\int_0^1 x dx\) 和 \(\int_0^1 x^2 dx\) 的大小。\end{exercise}
\begin{exercise}[3] 求 \(\int_0^2 |x-1| dx\) 的值。\end{exercise}
\begin{exercise}[3] 用几何意义求 \(\int_{-1}^1 (1-|x|) dx\)。\end{exercise}
\begin{exercise}[4] 用定义计算 \(\int_0^1 x^3 dx\)。\end{exercise}
\begin{exercise}[4] 求 \(\int_0^1 \sqrt{1-x^2} dx\)（几何意义：四分之一圆）。\end{exercise}
\begin{exercise}[4] 求 \(\int_{-2}^2 (x^3+1) dx\)。\end{exercise}
\begin{exercise}[4] 证明 \(\int_a^b f(x)dx = -\int_b^a f(x)dx\)。\end{exercise}
\begin{exercise}[4] 比较 \(\int_0^1 e^x dx\) 和 \(\int_0^1 e^{x^2} dx\) 的大小。\end{exercise}
\begin{exercise}[5] 用定义证明 \(\int_0^1 x^k dx = \frac{1}{k+1}\)。\end{exercise}
\begin{exercise}[5] 求 \(\int_0^{\pi/2} \sin x dx\) 用黎曼和极限。\end{exercise}
\begin{exercise}[5] 证明 Darboux 定理：上和与下和趋于同一极限。\end{exercise}
\begin{exercise}[5] 求 \(\int_{-1}^1 \frac{1}{1+x^2} dx\)（几何意义不容易直接看出，以后会算）。\end{exercise}
\begin{exercise}[6] 证明 Dirichlet 函数在 \([0,1]\) 上不可积。\end{exercise}
\begin{exercise}[6] 若 \(f\) 在 \([a,b]\) 上连续，证明 \(f\) 可积。\end{exercise}
\begin{exercise}[6] 证明 \(\int_0^1 \frac{\sin x}{x} dx\) 收敛。\end{exercise}
\begin{exercise}[7] 证明：若 \(f\) 在 \([0,1]\) 上可积，则 \(\lim_{n\to\infty}\frac1n\sum_{k=1}^n f(\frac{k}{n}) = \int_0^1 f(x)dx\)。\end{exercise}
